\section{Introduction}
    관성 측정 장치(Inertial Measurement Unit, IMU)는 3축 가속도계, 3축 자이로스코프, 그리고 선택적으로 3축 지자기계로 구성된 센서 모듈로, 물체의 자세(attitude)와 운동 상태를 추정하는 데 널리 활용된다. 소형화 및 저전력화가 가능한 MEMS 기반 IMU는 무인 항공기, 자율 주행 로봇, 위성 자세 제어 시스템(ADCS) 등 다양한 분야의 핵심 센서로 자리잡고 있다.\\

    그러나 IMU는 본질적으로 측정 노이즈, 바이어스 드리프트, 스케일 팩터 오차 등의 오차 특성을 가지며, 이를 적절히 모델링하고 보정하지 않으면 자세 추정 성능이 급격히 저하된다. 특히 자이로스코프의 각속도를 단순 적분하여 자세를 구하는 방식은 시간이 지남에 따라 누적 오차가 발생하는 드리프트(drift) 문제를 피할 수 없다. 이를 해결하기 위해 가속도계와 지자기계를 융합하는 센서 퓨전 기법이 필수적으로 요구된다.\\

    본 보고서는 MPU-9250 IMU를 대상으로 센서 노이즈 특성 분석, 캘리브레이션, 그리고 자세 추정 알고리즘의 설계 및 구현을 다룬다. 구체적인 목표는 다음과 같다.\\

    \begin{itemize}
        \item IMU 노이즈 성분의 확률적 모델링 및 Allan Variance, PSD 분석을 통한 센서 특성 정량화\cite{smith2023}
        \item 가속도계 정적 캘리브레이션 및 지자기계 Hard Iron / Soft Iron 보정
        \item Complementary Filter, Mahony/Madgwick Filter, Extended Kalman Filter(EKF)의 설계 및 성능 비교
        \item Python 시뮬레이션을 통한 알고리즘 검증 및 실제 하드웨어(Arduino + MPU-9250) 기반 실험 결과 분석
    \end{itemize}

    본 보고서의 구성은 다음과 같다. 2장에서는 IMU 센서의 기본 원리를 기반으로 Acclerometer와 Gyroscope, Magnetometer의 원리를 고찰함으로써 측정 모델을 정의한 이후, IMU 센서에서 흔히 논의되는 노이즈에 대해 논한다. 3장에서는 IMU 센서의 칼리브레이션 과정에 대해 각 Accelrometer, Gyroscope, Magnetometer에서의 칼리브레이션 과정을 정의한다. 이를 통해 IMU 센서의 정적 상태 추정 방법을 정의하고, 4장에서 IMU 센서의 좌표계를 정의한 이후 동적 상태 추정을 위한 Kinematic Model과 Measurement Model을 정의한다. 5장은 관성 항법 장치(Inertial Navigation System, INS)에서 센서 데이터를 정제하기 위한 AHRS 알고리즘들을 소개한다. 해당 AHRS 알고리즘들을 직접 구현함으로써 6장에서는 python을 통해 생성한 IMU 시뮬레이션 데이터를 기반한 상태 추정 알고리즘을 검증하고 이후 실제 IMU 센서(MPU9250)을 활용하여 센서 데이터 실측을 통해 앞서 설명한 AHRS 알고리즘의 성능을 비교분석한다. 7장은 결론 및 향후 과제를 제시함으로써 IMU 센서를 기반으로 제작한 AHRS 시스템의 상태 추정의 한계와 활용점을 논의한다.\\